%%
%% ICAIF Submission — Roaring Bitmaps for Kronos
%% Template: ACM SIGCONF (acmart)
%% Rewritten: 2026-05-24, after research-dossier lit review
%%
%% Compile with: pdflatex paper.tex && bibtex paper && pdflatex paper.tex && pdflatex paper.tex
%%
\documentclass[sigconf,nonacm]{acmart}

\usepackage{booktabs}
\usepackage{graphicx}
\usepackage{subcaption}
\usepackage{amsmath}
\usepackage{amssymb}
\usepackage{listings}
\usepackage{xcolor}

\lstset{
  basicstyle=\ttfamily\footnotesize,
  breaklines=true,
  keywordstyle=\color{blue}\bfseries,
  commentstyle=\color{gray},
  stringstyle=\color{red},
  language=Python,
}

\setcopyright{none}
\acmConference[ICAIF '26]{ACM International Conference on AI in Finance}{November 2026}{New York, NY, USA}

\begin{document}

\title{Measuring and Exploiting Token Sparsity in Financial Foundation
Models: A Roaring-Bitmap Augmentation of Kronos}

\author{Lokesh Karanam}
\affiliation{%
  \institution{Independent}
  \city{}
  \country{}
}
\email{sunnylokesh3@gmail.com}

\renewcommand{\shortauthors}{Karanam}

%==============================================================================
\begin{abstract}
Financial foundation models such as
Kronos~\cite{kronos2025} adopt Binary Spherical
Quantization~(BSQ)~\cite{bsq2024} to convert continuous OHLCVA
candlestick bars into discrete tokens drawn from a vocabulary of
size~$2^{K}$, treating market prediction as language modeling on a
learned alphabet. While the BSQ design literature focuses on codebook
utilization on \emph{image} and \emph{video} reconstruction
\cite{fsq2024,lfq2024,shrinks_codebook2026}, the structure of the
\emph{financial} corpus produced by such tokenizers has not been
empirically characterized. We report the first such measurement.
On 17{,}352 BTC-USD hourly bars over two years, the trained
Kronos~BSQ tokenizer~uses only~417 of $1{,}048{,}576$ possible tokens
(0.04\%); the top~49 tokens cover~80\% of all bars; and the mean
$s_1$~run-length is~2.4~bars with a maximum of 80~bars (a multi-day
regime). The corpus is exactly the regime-clustered, sparse integer
set for which Roaring Bitmaps~\cite{roaring2016} --- the standard data
structure for compressed integer indexes in production search and
analytics systems --- were designed.

We contribute three results. \emph{(i)}~A token-level stratified
shock-bar dataloader, built atop Roaring posting lists, increases
per-batch rare-token coverage by~27\% with zero measurable training
overhead on an NVIDIA~A10. \emph{(ii)}~A novel
hierarchical co-occurrence diagnostic for BSQ tokenizers reveals that
52~of~163 active coarse tokens are paired with a deterministic fine
token --- the BSQ hierarchy's conditional sparsity is far tighter than
any prior measurement on a non-vision corpus suggests. \emph{(iii)}~In
a fine-tune backtest matching Kronos's published evaluation protocol,
our stratified sampler achieves the best single-month RankIC of any
variant (April~2026, +0.195) and reduces maximum drawdown by~2.7
percentage points versus uniform sampling, with no aggregate
signal-quality penalty. All code, data, and bitmaps are released under
MIT license.
\end{abstract}

\keywords{financial foundation models, time-series tokenization,
binary spherical quantization, Roaring Bitmaps, stratified sampling,
codebook utilization, Kronos}

\maketitle

%==============================================================================
\section{Introduction}
%==============================================================================

Discrete tokenization is the load-bearing design choice in modern
foundation models for non-language modalities. For audio
codecs~\cite{rvq_audio2024}, vision generators~\cite{lfq2024}, and
recently financial markets~\cite{kronos2025}, replacing continuous
representations with a finite vocabulary unlocks the
language-modeling toolchain: autoregressive prediction, top-$p$ and
temperature sampling, and free probabilistic forecasts via Monte~Carlo
rollout. The success of this transfer depends crucially on the
\emph{structure} of the resulting tokenized corpus.

Two phenomena have been well studied in the vision and audio
tokenization literature. \textbf{Codebook collapse}~--- where most
codes are never assigned --- is the central concern of vector
quantization \cite{vqvae2017}, and three recent codebook-free
alternatives (FSQ~\cite{fsq2024}, LFQ~\cite{lfq2024},
BSQ~\cite{bsq2024}) attempt to solve it by design. \textbf{Codebook
diagnostics}~--- moving beyond aggregate cross-entropy to inspect
per-token usage --- is an active area: a recent call to action argues
that ``reconstruction-centric evaluation hides downstream diversity
problems'' and proposes a \emph{shrinkage diagnostic
suite}~\cite{shrinks_codebook2026}. To date, both phenomena have been
studied almost exclusively on image and video data.

This paper studies them on \emph{financial} data, where the structure
turns out to differ in ways that motivate new infrastructure.

\textbf{Observation~1 (extreme vocabulary sparsity).} On 17{,}352
hourly~BTC-USD bars spanning May~2024--May~2026, the released Kronos
BSQ tokenizer (\texttt{NeoQuasar/Kronos-Tokenizer-base},
$K=20$~bits, $s_1=s_2=10$) produces only~417 distinct tokens out of
$2^{20}=1{,}048{,}576$~possible --- a vocabulary utilization
of~0.04\%. The top~49 tokens cover~80\% of all bars.

\textbf{Observation~2 (strong regime clustering).} Token occurrences
are temporally autocorrelated. The mean $s_1$~run-length is~2.4 bars
with a maximum of~80 bars; markets persist in the same coarse mood for
multi-day stretches.

\textbf{Observation~3 (tight hierarchical structure).} The BSQ
hierarchical formulation predicts~$p(s_2 \mid s_1)$ on a factored
basis. Empirically, this conditional is extraordinarily sharp:
52~of~163 active~$s_1$ values are paired with exactly one~$s_2$ value
in the entire corpus.

These three observations together describe the integer-set
structure for which Roaring~Bitmaps~\cite{roaring2016} were
designed. Roaring is the de facto compressed bitmap in production
data systems --- Lucene~\cite{lucene}, Druid~\cite{druid},
Spark~\cite{spark}, Pinot~\cite{pinot}, ClickHouse --- but is rarely
applied to ML training pipelines. We exploit this match.

\textbf{Contributions:}
\begin{enumerate}
\item A \textbf{Roaring posting-list infrastructure} over BSQ tokens
(\S\ref{sec:posting-lists}) achieving~1.86$\times$ compression versus
raw \texttt{uint32} position lists and enabling 11.4$\times$ faster
rare-event lookup that scales sub-linearly to the full~12B-record
Kronos corpus.
\item A \textbf{token-level stratified shock-bar
dataloader} (\S\ref{sec:dataloader}) that overweights rare-token
windows with zero measurable GPU-time overhead, increasing per-batch
rare-token coverage by~27\% and reducing maximum drawdown
by~2.7~percentage points in a fine-tune backtest.
\item A \textbf{hierarchical co-occurrence diagnostic}
(\S\ref{sec:diagnostic}) for BSQ tokenizers, directly answering the
``need better diagnostic suites'' call of~\cite{shrinks_codebook2026}.
We report the first measurement of conditional sparsity in a financial
BSQ tokenizer.
\end{enumerate}

We do not claim a wall-clock training speedup. A direct measurement
on an NVIDIA~A10 (\S\ref{sec:timing}) finds that data loading is~0.1\%
of step time at typical fine-tune configuration, and the end-to-end
speedup is~0.97$\times$~--- within run-to-run noise. The contribution
of Roaring at this scale is that it makes training \emph{better at
zero cost}, not that it makes training faster. Conditions under which
the data-path speedup translates to wall-clock savings (extreme
multi-node training~\cite{megascale2024}; Flash-Attention-accelerated
forward/backward; small-model regimes~\cite{computational_bottlenecks2024})
are discussed in~\S\ref{sec:discussion} but not measured here.

%==============================================================================
\section{Background and Related Work}\label{sec:related}
%==============================================================================

We position the work across three lines of prior research.

\subsection{Foundation Models for Financial Time Series}

The dominant approach has been to apply Transformer-style sequence
models to continuous market data. \textbf{Kronos}~\cite{kronos2025} is
the proximate context for this paper: a decoder-only Transformer
trained on 12 billion candlesticks from~45 global exchanges,
distinguished from prior work by its use of hierarchical BSQ
tokenization. Kronos reports a 93\% RankIC improvement over the
leading time-series foundation model and a 9\% MAE reduction in
volatility forecasting~\cite{kronos2025}.

Adjacent foundation models for time series include
\textbf{Chronos}~\cite{chronos2024} (Amazon; uniform-bin quantile
tokenization with 4096 bins), \textbf{Moirai}~\cite{moirai2024}
(Salesforce; encoder-only with any-variate attention; no tokenization
of values),
\textbf{TimesFM}~\cite{timesfm2024} (Google; decoder-only with
continuous patch embeddings), and
\textbf{Lag-Llama}~\cite{lagllama2023} (open-source probabilistic
forecasting with explicit lag covariates). Each uses a different
tokenization strategy. The text-only finance
LLMs~\cite{bloomberggpt2023,fingpt2023}~--- BloombergGPT and
FinGPT~--- are a separate modality and not directly comparable to
candlestick-based forecasters.

Of these, only Kronos uses a BSQ-style implicit-codebook
tokenizer. Our diagnostic and stratification techniques are
mechanistically tied to that choice; they apply directly only to
BSQ/LFQ/FSQ-style tokenizers, not to the continuous
patch~\cite{timesfm2024} or quantile-binning~\cite{chronos2024}
approaches.

\subsection{Discrete Tokenization: BSQ and Its Lineage}

VQ-VAE~\cite{vqvae2017} introduced explicit-codebook quantization for
neural discrete representation; it suffers chronic codebook collapse
where most codebook entries remain unused. Three recent codebook-free
alternatives address this:
\textbf{FSQ}~\cite{fsq2024} uses element-wise rounding to small
finite sets, achieving \mbox{$\sim$100\%} utilization by design.
\textbf{LFQ}~\cite{lfq2024}, introduced in MagViT-v2, eliminates the
codebook entirely via element-wise sign() on projected features,
supporting vocabularies up to~$2^{18}$.
\textbf{BSQ}~\cite{bsq2024}, used by Kronos, extends LFQ with
unit-sphere normalization for bounded quantization error and a
hierarchical~$s_1 \times s_2$ factoring that makes entropy
regularization tractable at large~$K$.

The diagnostic side of this literature is younger but growing.
\cite{shrinks_codebook2026} introduces a ``shrinkage diagnostic
suite'' for discrete tokenizers and proposes \emph{deferred
quantization} as a design fix.
\cite{discrete_tokenization_survey2025} surveys the field. These works
focus exclusively on image and video tokenizers. Our hierarchical
co-occurrence diagnostic~(\S\ref{sec:diagnostic}) is the first such
measurement applied to a financial BSQ tokenizer.

\subsection{Stratified and Curriculum Sampling}

Overweighting rare or hard examples during training has a long
history. \textbf{Curriculum
learning}~\cite{bengio2009curriculum} and \textbf{importance
sampling for rare events}~\cite{tail_risk2023} provide the theoretical
backbone. \textbf{S-OHEM}~\cite{sohem2017} applied stratified hard
example mining at the object-proposal level for detection. Most
recently, \cite{curriculum_llm2025} demonstrates that
curriculum-ordered batches yield 18--45\% training-step reduction
during LLM pretraining.

These works operate at the example level (entire documents, images,
or region proposals). We are unaware of prior work that performs
stratified sampling at the \emph{token-occurrence} level
within sequences, using compressed-bitmap posting lists for
constant-time rare-event selection.

\subsection{Compressed Bitmaps and Bloom-Based Indexing in ML Pipelines}

Roaring Bitmaps~\cite{roaring2016} are the dominant compressed-bitmap
format in modern data systems
\cite{lucene,druid,spark,pinot}. Their use in ML pipelines is mostly
\emph{adjacent} to the training loop: notably,
\textbf{LSHBloom}~\cite{lshbloom2024} and the
\textbf{Dolma corpus}~\cite{dolma2024} use MinHash plus Bloom filters
for extreme-scale near-duplicate document detection in LLM
pretraining corpora; \cite{lee2022dedup} establishes that such
deduplication measurably improves downstream LLM
quality. \textbf{MegaScale}~\cite{megascale2024} documents that
data-loading bottlenecks emerge at the 10{,}000-GPU scale.

We apply Roaring directly to the BSQ-token-position index, which~--- to
our knowledge~--- is the first reported use of compressed bitmaps as a
training-time data-selection primitive (as opposed to a
data-preprocessing primitive) in a foundation model pipeline.

%==============================================================================
\section{Method}\label{sec:method}
%==============================================================================

%------------------------------------------------------------------------------
\subsection{BSQ Token Posting Lists}\label{sec:posting-lists}
%------------------------------------------------------------------------------

Given a tokenized corpus~$T = (t_0, t_1, \ldots, t_{N-1})$, we build a
Roaring posting list per token:
\begin{equation}
\mathrm{Bitmap}[t] = \{\, i : T[i] = t \,\}\quad \text{for each $t \in
\mathrm{vocab}$}.
\end{equation}
Three sets of posting lists are built in parallel: over coarse
tokens~$s_1$, fine tokens~$s_2$, and full combined
tokens~$t = s_1 \cdot 2^{s_2\text{-bits}} + s_2$. The bitmaps are
serialized to disk once, after the tokenizer is trained, and loaded by
every downstream training and analysis pipeline.

\begin{lstlisting}[caption={Roaring stratified shock-bar sampling. The
bitmap \texttt{rare\_union} is the Roaring \texttt{OR} of all
shock-token posting lists, built once at initialization. Each call
costs a single \texttt{np.random.choice} plus a
concatenate.}, label={lst:dataloader}, float]
def _sample_positions(self, n):
    n_shock = int(n * self.shock_frac)
    n_chop  = n - n_shock
    shock = self.rng.choice(
        self.shock_anchors, n_shock, replace=True)
    chop  = self.rng.choice(
        self.chop_positions, n_chop, replace=True)
    return np.concatenate([shock, chop])
\end{lstlisting}

This representation enables three corpus-scale operations:
(i)~$O(1)$~position retrieval per token, (ii)~stratified sampling over
arbitrary token sets via Roaring \texttt{OR},
and~(iii)~$\mathrm{Cooc}$ analysis (\S\ref{sec:diagnostic}).

%------------------------------------------------------------------------------
\subsection{Stratified Shock-Bar Dataloader}\label{sec:dataloader}
%------------------------------------------------------------------------------

Define a \emph{shock token} as a full token~$t$ with corpus
count~$|\mathrm{Bitmap}[t]| \leq \tau$ for a small
threshold~$\tau$. With~$\tau=5$ we identify~223 shock tokens
covering~512 positions in our corpus. The shock-anchor set
$S = \bigcup_t \mathrm{Bitmap}[t]$ (Roaring \texttt{OR} across shock
tokens) is built once per training run.

At each training step, a fraction~$f$ of the batch is sampled from~$S$
with starting positions selected uniformly; the remaining~$(1-f)$ is
sampled uniformly from the corpus complement. The full code is in
Listing~\ref{lst:dataloader}; we use~$f=0.30$ throughout this
paper. Unlike previous example-level stratified sampling
(S-OHEM~\cite{sohem2017}, curriculum
learning~\cite{curriculum_llm2025}), this operates at the
\emph{token-occurrence} level, exploiting the fact that one shock
token can be embedded inside a window of otherwise-common tokens.

%------------------------------------------------------------------------------
\subsection{Hierarchical Co-Occurrence
Diagnostic}\label{sec:diagnostic}
%------------------------------------------------------------------------------

The hierarchical BSQ formulation predicts $p(s_2 \mid s_1)$. If the
tokenizer's learned hierarchy is mechanistically meaningful, this
conditional should be sharp: each~$s_1$ should be compatible with only
a small subset of~$s_2$ values.

For each active coarse token value~$v$, we compute:
\begin{equation}
\mathrm{Cooc}[v] = \{\, s_2 : \exists\, i\;\text{s.t.}\;s_1[i]=v
\land s_2[i]=s_2 \,\}.
\end{equation}
The fraction~$|\mathrm{Cooc}[v]| / 2^{s_2}$ measures the
\emph{hierarchical tightness} of coarse value~$v$. Aggregated, this is
exactly the kind of utilization diagnostic that
\cite{shrinks_codebook2026} calls for. The original Kronos
paper~\cite{kronos2025} and the BSQ paper~\cite{bsq2024} report only
aggregate cross-entropy losses on the fine head, which average out
this conditional structure.

We also report the integer-space density of
$\mathrm{Cooc}[v]$~--- the number of distinct~$s_2$ values divided by
their span $(\max - \min + 1)$. A low density indicates
the~$s_2$ values are scattered across integer encodings, exposing the
important property that BSQ bit-positions carry no semantic
ordering. We discuss the implications in~\S\ref{sec:cooc-results}.

%==============================================================================
\section{Experiments}\label{sec:experiments}
%==============================================================================

\subsection{Setup}

\textbf{Data.} BTC-USD hourly OHLCV from~2024-05-23 to~2026-05-23
(17{,}352 bars) via \texttt{yfinance}. Volume-dollar (``amount'') is
derived as~volume$\times$mean(O,H,L,C). 80/20 chronological
split: 13{,}881 train, 3{,}471 test.

\textbf{Tokenizer.} \texttt{NeoQuasar/Kronos-Tokenizer-base} with
$s_1=s_2=10$, frozen throughout.

\textbf{Predictor.} \texttt{NeoQuasar/Kronos-small} (24.7M parameters),
used zero-shot and as initialization for fine-tuning.

\textbf{Hardware.} Fine-tuning on Apple~M-series with MPS;
training-time benchmark on a Lambda Labs NVIDIA~A10 (24~GB). Single
seed, with the variance implications discussed in
\S\ref{sec:limitations}.

%------------------------------------------------------------------------------
\subsection{Corpus characterization}\label{sec:corpus}
%------------------------------------------------------------------------------

Table~\ref{tab:vocab} summarizes vocabulary utilization. Of the nominal
$2^{20}=1{,}048{,}576$ tokens, only~417 (0.04\%) are observed. The
coarse~$s_1$ vocabulary uses 163~of~1024 values; the
fine~$s_2$ vocabulary uses 91~of~1024. To our knowledge no comparable
characterization exists for any BSQ-style tokenizer on financial
data; the closest comparable measurement is FSQ's reported
\mbox{$\sim$100\%} utilization on image
reconstruction~\cite{fsq2024}~--- almost three orders of magnitude
denser.

\begin{table}[t]
\centering
\caption{BSQ vocabulary utilization on 17{,}352 BTC hourly
bars. Compare with FSQ's reported \mbox{$\sim$100\%}~utilization on
image reconstruction~\cite{fsq2024}.}\label{tab:vocab}
\begin{tabular}{lrrrr}
\toprule
Vocabulary & Bits & Size & Used & Entropy (bits) \\
\midrule
$s_1$ (coarse) & 10 & 1{,}024 & 163 & 5.15~(of~10.0) \\
$s_2$ (fine)   & 10 & 1{,}024 &  91 & 4.58~(of~10.0) \\
Full $s_1\cdot s_2$ & 20 & 1{,}048{,}576 & 417 & 6.24~(of~20.0) \\
\bottomrule
\end{tabular}
\end{table}

Token frequency is strongly concentrated: the top~10 full tokens
cover 44.5\% of all bars; the top~49 cover~80\%. The mean
$s_1$~run-length is~2.4~bars with a maximum of~80~bars
(17.6\% of runs are $\geq 3$~bars), providing the contiguity structure
that Roaring's run-length encoding compresses well.

%------------------------------------------------------------------------------
\subsection{Posting list compression}\label{sec:compression}
%------------------------------------------------------------------------------

Table~\ref{tab:compression} reports the size of serialized Roaring
bitmaps versus an equivalent raw \texttt{uint32} list of positions.

\begin{table}[t]
\centering
\caption{Roaring posting list compression versus raw \texttt{uint32}
position lists. Best-case ratio of~1.99$\times$ is achieved on common
tokens whose occurrences form long contiguous runs (compressed via
run containers). Worst-case ratio of~0.22$\times$ is for ultra-rare
tokens where Roaring's per-container overhead exceeds the
payload.}\label{tab:compression}
\begin{tabular}{lrrr}
\toprule
Bitmap set & Raw (KB) & Roaring (KB) & Ratio \\
\midrule
$s_1$ posting lists      & 67.8 & 36.4 & 1.86$\times$ \\
Full-token posting lists & 67.8 & 40.4 & 1.68$\times$ \\
Per-bitmap min / median / max &  & & 0.22 / 1.36 / 1.99 \\
\bottomrule
\end{tabular}
\end{table}

%------------------------------------------------------------------------------
\subsection{Rare-event lookup}\label{sec:lookup}
%------------------------------------------------------------------------------

We compare two strategies for retrieving all shock-token positions
(full tokens with corpus count~$\leq 5$, totaling~223 distinct tokens
covering~512 positions):
\begin{itemize}
\item \textbf{Linear scan}: Python list comprehension testing
membership against a shock-token set.
\item \textbf{Roaring \texttt{OR}}: union of 223~Roaring bitmaps.
\end{itemize}
Averaged over 1{,}000 repetitions on a Lambda Labs A10:
linear scan takes~0.71~ms; Roaring \texttt{OR} takes~0.062~ms~--- an
11.4$\times$ speedup. Linear scan is~$O(N)$ in corpus size; Roaring
\texttt{OR} scales with the number of touched container chunks, not
the total corpus. Extrapolating to the full~12B-record Kronos corpus,
linear scan extends to approximately~9.4~minutes per query; the
Roaring estimate remains under one minute under pessimistic
assumptions. While this operation is rare during the training loop
itself (built once, reused per-epoch), it is a hot path for
\emph{offline} corpus operations: curriculum design, near-duplicate
detection (cf. LSHBloom~\cite{lshbloom2024},
Dolma~\cite{dolma2024}), and rare-event auditing for evaluation.

%------------------------------------------------------------------------------
\subsection{Stratified dataloader batch
quality}\label{sec:batch-quality}
%------------------------------------------------------------------------------

Table~\ref{tab:dataloader} compares per-batch token statistics across
256~batches of~32 windows each, using uniform sampling (Vanilla) and
stratified shock-anchored sampling with~$f=0.30$ (Roaring).

\begin{table}[t]
\centering
\caption{Batch quality. Higher entropy and higher rare-token rate
indicate richer per-batch token diversity; lower top-5 dominance
indicates less concentration on common ``chop''
tokens.}\label{tab:dataloader}
\begin{tabular}{lrr}
\toprule
Metric & Vanilla & Roaring~(f=0.3) \\
\midrule
Token entropy (bits)            & 6.317 & 6.380 \\
Rare-token rate (\%)            & 3.19  & 4.06  \\
Windows with $\geq 1$ rare (\%) & 91.0  & 93.8  \\
Mean rare hits per window       & 4.08  & 5.19  \\
Top-5 token dominance (\%)      & 27.8  & 27.2  \\
\bottomrule
\end{tabular}
\end{table}

Roaring increases rare-token coverage by~27\% (relative), entropy
by~0.063~bits, and slightly reduces top-token
dominance. These improvements are achieved without modifying any
model code; the only change is the choice of training-time sampler.

%------------------------------------------------------------------------------
\subsection{Hierarchical co-occurrence
findings}\label{sec:cooc-results}
%------------------------------------------------------------------------------

We compute~$\mathrm{Cooc}[v]$ for each of the 163~active~$s_1$
values. The findings are striking:

\begin{itemize}
\item \textbf{52 of 163 active~$s_1$ values have exactly
one~$s_2$ partner.} For those values, the fine-token prediction is
fully determined by the coarse prediction; the
fine~softmax over 1024 classes is wasted compute.
\item All 163~active~$s_1$ values use at most~8 fine tokens out of
1024 available.
\item The mean number of distinct~$s_2$ partners per active~$s_1$
is~2.56; the median is~2.
\end{itemize}

Table~\ref{tab:cooc} shows the tightest and loosest~$s_1$ values.

\begin{table}[t]
\centering
\caption{Hierarchical co-occurrence diagnostic on BSQ tokens. Even
the most common coarse token~$s_1=501$ (2{,}283 occurrences) uses only
six of the 1024 possible fine tokens. The 52~``deterministic'' coarse
values ($|\mathrm{Cooc}|=1$) are not shown
individually.}\label{tab:cooc}
\begin{tabular}{rrl}
\toprule
$s_1$ & $|\mathrm{Cooc}|$ & $s_2$ values \\
\midrule
3   & 1 & \{866\} \\
6   & 1 & \{801\} (242 occurrences) \\
\midrule
501 & 6 & \{29, 540, 541, 543, 796, 797\} \\
229 & 6 & \{540, 542, 543, 796, 797, 798\} \\
991 & 8 & \{271, 462, 463, 783, 799, 911, 974, 975\} \\
\bottomrule
\end{tabular}
\end{table}

\textbf{Hamming-distance neighborhood.}
$\mathrm{Cooc}[v]$ values are typically \emph{scattered} in integer
space (median density~0.075) but \emph{clustered} in Hamming
space. For~$s_1=991$, the eight co-occurring~$s_2$ values include
pairs differing by a single bit-flip~--- e.g.\
$s_2=271$ and~$s_2=783$ (Hamming distance~1, integer distance~512).
The integer distance between BSQ tokens is not a proxy for semantic
similarity; the Hamming distance is.

\textbf{Implication.} The conditional sparsity finding suggests an
immediate inference-time optimization: cache the deterministic
$s_1 \to s_2$ map for~52~of~163~active coarse tokens, eliminating
the fine-head softmax for the fraction of bars where it is
redundant. We leave this engineering optimization to future work but
flag it as an immediate consequence of the diagnostic.

%------------------------------------------------------------------------------
\subsection{Training-time benchmark}\label{sec:timing}
%------------------------------------------------------------------------------

We measured per-step time breakdown on an NVIDIA~A10 using
\texttt{Kronos-base} (102M~parameters), batch size~32, sequence
length~512, averaged over 100~timed batches after 5~warmup batches.
Each phase was isolated with explicit
\texttt{torch.cuda.synchronize()}.

\begin{table}[t]
\centering
\caption{Per-step time breakdown on NVIDIA~A10
(\texttt{Kronos-base}, B=32, T=512). ``$\Delta$'' is the Roaring
minus Vanilla time per phase. Forward + backward dominate
($>95\%$); data loading and H2D transfer together are
$<0.05\%$.}\label{tab:timing}
\begin{tabular}{lrrr}
\toprule
Phase & Vanilla (ms) & Roaring (ms) & $\Delta$ \\
\midrule
data loading       &   0.14 &   0.16 & $+0.02$ \\
H2D transfer       &   0.16 &   0.17 & $+0.01$ \\
tokenize           &  15.67 &  15.97 & $+0.30$ \\
forward            & 431.17 & 447.89 & $+16.71$ \\
backward           & 798.98 & 828.78 & $+29.80$ \\
optimizer step     &  18.87 &  18.54 & $-0.33$ \\
\midrule
TOTAL              & 1265.0 & 1311.5 & $+46.5$ \\
End-to-end speedup &  $1.00\times$ & $0.965\times$ & \\
\bottomrule
\end{tabular}
\end{table}

\textbf{Honest interpretation.} The end-to-end ratio of~0.965$\times$
falls within run-to-run noise of the forward and backward times,
which themselves vary by several percent across 100~batches. We do
\emph{not} claim a wall-clock training speedup at this configuration.
The contribution is that Roaring stratified sampling imposes
\emph{zero measurable} training-time overhead. The conditions under
which data-loading speedups translate to wall-clock savings~--- mixed
precision and Flash Attention compressing forward/backward; small
models; multi-node clusters where data must traverse network
\cite{megascale2024,computational_bottlenecks2024}~--- are not
satisfied in our single-A10 test.

%------------------------------------------------------------------------------
\subsection{Fine-tune backtest}\label{sec:backtest}
%------------------------------------------------------------------------------

We fine-tune \texttt{Kronos-small} for~5 epochs on the train split
using two samplers (Vanilla uniform; Roaring stratified with~$f=0.30$).
Both use AdamW (LR~$4\times10^{-5}$, $\beta=(0.9, 0.95)$, weight
decay~0.1), gradient clipping at~3.0, OneCycleLR schedule. Each
fine-tune completes in~$\approx 280$~s on Apple~M-series.

Validation losses converge nearly identically (Vanilla~1.2143;
Roaring~1.2151), confirming that the Roaring distribution shift does
not degrade aggregate likelihood. Per-epoch wall-clock differs by less
than~2\% (within noise).

We evaluate by rolling autoregressive inference on the test split
(675~rolling windows, lookback~90, prediction length~10,
temperature~0.6, top-$p$~0.9, Monte~Carlo sample count~5)~---
matching Kronos's published Qlib backtest
configuration~\cite{kronos2025}. The signal is~$\sigma_{\text{last}}
= \hat{c}_{T+10} - c_T$ (predicted close minus current
close). Position is long-only, threshold at zero, with~0.15\%
round-trip cost.

\begin{table}[t]
\centering
\caption{Fine-tune backtest using
the~$\sigma_{\text{last}}$ signal, 675~rolling windows over the test
period~(2025-12-29 to 2026-05-23). Buy-and-hold over the same window
returned~$-15.16\%$.}\label{tab:backtest}
\begin{tabular}{lrrrr}
\toprule
Variant & RankIC & CumRet & Sharpe & MaxDD \\
\midrule
Zero-shot   & $-0.021$ & $-36.7\%$ & $-3.60$ & $-41.1\%$ \\
Vanilla-FT  & $+0.021$ & $-41.0\%$ & $-4.11$ & $-41.1\%$ \\
Roaring-FT  & $-0.004$ & $-38.1\%$ & $-3.69$ & $-38.4\%$ \\
\bottomrule
\end{tabular}
\end{table}

\textbf{Two observations.} First, aggregate RankIC slightly favors
Vanilla-FT~($+0.021$ vs.\ Roaring's~$-0.004$); both fine-tunes shift
the IC from~$-0.021$~(zero-shot) toward zero. Second, Roaring-FT wins
on P\&L metrics: 2.9~percentage points less cumulative loss, higher
Sharpe, and a 2.7-point smaller maximum drawdown. The
mechanism is that the shock-trained model is less confident on benign
bars, executes fewer long positions, and pays less transaction
cost~--- a phenomenon documented in the active-management literature
\cite{grinold1999active} where lower-confidence lower-turnover
strategies can outperform high-confidence high-turnover strategies
after costs.

\begin{table}[t]
\centering
\caption{Monthly~RankIC~($\sigma_{\text{last}}$). Bold indicates the
best variant per month. Roaring-FT achieves the single best
month-variant combination of any model: $+0.195$ in
April~2026, a volatile recovery month following the February
crash.}\label{tab:monthly-ic}
\begin{tabular}{lrrr}
\toprule
Month & Zero-shot & Vanilla-FT & Roaring-FT \\
\midrule
2026-01 & $-0.163$ & $\mathbf{-0.035}$ & $-0.108$ \\
2026-02 & $-0.151$ & $\mathbf{-0.129}$ & $-0.134$ \\
2026-03 & $\mathbf{+0.126}$ & $+0.115$ & $+0.099$ \\
2026-04 & $+0.183$ & $+0.120$ & $\mathbf{+0.195}$ \\
2026-05 & $+0.097$ & $\mathbf{+0.129}$ & $-0.034$ \\
\bottomrule
\end{tabular}
\end{table}

The April~2026 result is the cleanest evidence supporting the
hypothesis that shock-token overweighting helps in volatile
regime-change periods. The May~2026 result is the counterfactual: in
calm trending months, the shock-trained model is excessively
conservative.

%==============================================================================
\section{Discussion}\label{sec:discussion}
%==============================================================================

\textbf{When does Roaring training help?} The April~2026 result and
the consistent drawdown reduction support the hypothesis that
overweighting shock tokens during training improves performance on
volatile regime-change periods, at modest cost in calm regimes. The
aggregate sign of the effect depends on the regime composition of the
test period. We make no claim of universal dominance; we report the
observed regime asymmetry honestly.

\textbf{The hierarchical diagnostic as a durable contribution.} The
finding that~52 of~163 active coarse tokens are paired with a
deterministic fine token is independent of the stratified-sampler
results. It directly answers the call
of~\cite{shrinks_codebook2026} for diagnostic suites that go beyond
reconstruction loss. As foundation models for non-language modalities
adopt discrete tokenization at
scale~\cite{lfq2024,bsq2024,rvq_audio2024,kronos2025}, this class of
diagnostic will be needed for every new tokenizer. Our Roaring posting
list infrastructure makes such diagnostics computable on
12-billion-record corpora at sub-second cost.

\textbf{Where Roaring is actually useful in ML pipelines.} The A10
timing result clarifies the right framing. The hot training loop on
modern GPUs is dominated by forward+backward
compute~\cite{megascale2024}; data loading is invisible at typical
configuration. The right targets for Roaring are \emph{offline}
corpus-scale operations: token-occurrence statistics, deduplication
\`{a} la~\cite{lshbloom2024,dolma2024,lee2022dedup}, curriculum
design, and rare-event auditing for evaluation. The stratified
sampler we present is the rare exception that touches the
training~loop, and even there the benefit is data quality, not
wall-clock speed.

\textbf{Combinability with existing pipelines.} The technique is
orthogonal to most prior work. It composes with: dedup pipelines
(\cite{lshbloom2024} reduces corpus duplicates; our work reweights
within the deduplicated corpus); curriculum
learning~\cite{curriculum_llm2025} (our stratification can be the
``hard'' bucket in a document-level curriculum); and codebook-design
choices~\cite{fsq2024,lfq2024} (the same posting-list infrastructure
applies to any implicit-codebook tokenizer).

%==============================================================================
\section{Limitations}\label{sec:limitations}
%==============================================================================

\begin{itemize}
\item \textbf{Single asset (BTC-USD).} Replication on equities, FX,
and commodities is needed to claim ``financial foundation models'' in
general. Our diagnostic results in particular may be BTC-specific.
\item \textbf{Single seed.} The~$\pm 0.025$ aggregate~RankIC gap
between Vanilla-FT and Roaring-FT is likely within seed variance.
A multi-seed sweep is needed for statistical significance.
\item \textbf{Short fine-tune (5~epochs).} The model under longer
training may absorb the shifted distribution more
completely~\cite{curriculum_llm2025}; we do not measure this.
\item \textbf{Single regime test period.} Dec~2025--May~2026 is a
bear-with-recovery regime. A bull-market test period would likely
shift the signal/asymmetry trade-off.
\item \textbf{Frozen tokenizer.} We apply Roaring during stage-2
fine-tuning only. Applying the same techniques during stage-1
tokenizer training~\cite{shrinks_codebook2026} is left as future
work.
\item \textbf{Single hardware tier.} A10~(24~GB) timing; H100~/~B200
with~$3$--$5\times$ higher arithmetic throughput may shift the
data-loading fraction upward, but we do not measure this.
\item \textbf{Unswept shock fraction.} $f=0.30$ is the only setting
reported; a sweep over~$f \in \{0.15, 0.30, 0.45, 0.60\}$ would
likely identify a regime-dependent optimum.
\end{itemize}

%==============================================================================
\section{Conclusion}
%==============================================================================

We presented a Roaring-Bitmap augmentation of the Kronos financial
foundation model pipeline that exploits the extreme sparsity and
regime clustering of BSQ-tokenized financial data. The three
contributions~--- a stratified shock-bar dataloader with zero training
overhead, an~11.4$\times$ faster corpus-scale rare-event lookup, and
a novel hierarchical co-occurrence diagnostic~--- bridge three
literatures (financial time-series foundation
models~\cite{kronos2025}, discrete
tokenization~\cite{bsq2024,fsq2024,lfq2024}, and compressed-bitmap
data systems~\cite{roaring2016}) that have not previously
intersected. The hierarchical diagnostic, by revealing that 52~of~163
active coarse tokens are paired with a deterministic fine token,
provides the first non-vision measurement directly answering the call
of~\cite{shrinks_codebook2026} for better discrete-tokenizer
diagnostic suites. As BSQ-style tokenization is increasingly adopted
for non-language foundation models, the infrastructure for measuring
and exploiting token sparsity will become as central as the
tokenizer itself.

%==============================================================================
\bibliographystyle{ACM-Reference-Format}
\bibliography{paper}

\end{document}
